\documentclass[a4paper,11pt,reqno]{amsart}

\usepackage[utf8]{inputenc}
\usepackage[foot]{amsaddr}
\usepackage{amsmath,amsfonts,amssymb,amsthm,mathrsfs,bm}
\usepackage[margin=0.95in]{geometry}
\usepackage{color}
\usepackage[dvipsnames]{xcolor}

\input{toc-config.tex}

\usepackage{mathtools,enumerate,graphicx}
\usepackage{epstopdf}
\usepackage{hyperref}

\usepackage{latexsym}

\usepackage{pgfplots}
\pgfplotsset{compat=1.18}

\definecolor{CommentGreen}{rgb}{0.0,0.4,0.0}
\definecolor{Background}{rgb}{0.9,1.0,0.85}
\definecolor{lrow}{rgb}{0.914,0.918,0.922}
\definecolor{drow}{rgb}{0.725,0.745,0.769}

\usepackage{listings}
\usepackage{textcomp}
\lstloadlanguages{Matlab}%
\lstset{
    language=Matlab,
    upquote=true, frame=single,
    basicstyle=\small\ttfamily,
    backgroundcolor=\color{Background},
    keywordstyle=[1]\color{blue}\bfseries,
    keywordstyle=[2]\color{purple},
    keywordstyle=[3]\color{black}\bfseries,
    identifierstyle=,
    commentstyle=\usefont{T1}{pcr}{m}{sl}\color{CommentGreen}\small,
    stringstyle=\color{purple},
    showstringspaces=false, tabsize=5,
    morekeywords={properties,methods,classdef},
    morekeywords=[2]{handle},
    morecomment=[l][\color{blue}]{...},
    numbers=none, firstnumber=1,
    numberstyle=\tiny\color{blue},
    stepnumber=1, xleftmargin=10pt, xrightmargin=10pt
}

\numberwithin{equation}{section}
\synctex=1

\hypersetup{
    unicode=false, pdftoolbar=true, 
    pdfmenubar=true, pdffitwindow=false, pdfstartview={FitH}, 
    pdftitle={ELE2024 Coursework}, pdfauthor={A. Author},
    pdfsubject={ELE2024 coursework}, pdfcreator={A. Author},
    pdfproducer={ELE2024}, pdfnewwindow=true,
    colorlinks=true, linkcolor=red,
    citecolor=blue, filecolor=magenta, urlcolor=cyan
}


% CUSTOM COMMANDS
\renewcommand{\Re}{\mathbf{re}}
\renewcommand{\Im}{\mathbf{im}}
\newcommand{\R}{\mathbb{R}}
\newcommand{\N}{\mathbb{N}}
\newcommand{\C}{\mathbb{C}}
\newcommand{\lap}{\mathscr{L}}
\newcommand{\dd}{\mathrm{d}}
\newcommand{\smallmat}[1]{\left[ \begin{smallmatrix}#1 \end{smallmatrix} \right]}

%opening
\title[MATH 2860 (Elementary Differential Equations)]{Homework 7 for MATH 2860 (Elementary Differential Equations)}

\author[Emmanuel Atindama]{E. A. Atindama, PhD Mathematics}

\address[E. A. Atindama]{. Email addresses: \href{emmanuel.atindama@utoledo.edu}{emmanuel.atindama@utoledo.edu} 
% and 
% \href{mailto:a.student@qub.ac.uk}{a.student@qub.ac.uk}.
}
\thanks{Version 0.0.1. Last updated:~\today.}

\begin{document}
\maketitle

Due at 10:00 EST in class


\subsection*{Question Q1}
The complex-valued function \(f(t)\) depends on real variable \(t\):
\[
f(t)=\exp\big[(-1+i2\pi)t\big].
\]
\begin{enumerate}[(a)]
  \item Compute \(\Re f(t)\) and \(\Im f(t)\).
  \item Sketch \(\Re f(t)\) and \(\Im f(t)\) for \(0\le t<3\) (describe key features).
  \item For which \(t\) is \(f(t)\) purely real / purely imaginary?
\end{enumerate}

\begin{center}\setlength{\fboxsep}{8pt}\fcolorbox{yellow!20}{yellow!20}{\parbox{0.95\linewidth}{
\textbf{Solution}

\begin{enumerate}[(a)]
  \item Write \(f(t)=e^{-t}e^{i2\pi t}=e^{-t}\big(\cos(2\pi t)+i\sin(2\pi t)\big)\). Hence
    \[
    \Re f(t)=e^{-t}\cos(2\pi t),\qquad \Im f(t)=e^{-t}\sin(2\pi t).
    \]

  \item Key features:
    \begin{itemize}
      \item Both are sinusoidal with period \(T=1\). They decay from value \(1\) at \(t=0\) because of the product with \(e^{-t}\).
      \item Real part zeros: \[\cos(2\pi t)=0\Rightarrow t=\tfrac14 + \tfrac{n}{2},\ n\in\mathbb{Z}.\]
            In \([0,3)\): \(t=\tfrac14, \tfrac34, \tfrac54, \tfrac{3}{2}, \tfrac{7}{4}, \tfrac{9}{4}, \tfrac{11}{4}\).
      \item Imaginary part zeros: \[\sin(2\pi t)=0\Rightarrow t=\tfrac{n}{2},\ n\in\mathbb{Z}.\]
            In \([0,3)\): \(t=0,\tfrac12,1,\tfrac32,2,\tfrac52\).
      \item At \(t=0\): \(\Re f(0)=1,\ \Im f(0)=0\). Both decay in amplitude as \(e^{-t}\).
    \end{itemize}

    \begin{tikzpicture}
      \begin{axis}[
        width=12cm, height=7cm,
        xmin=0, xmax=3,
        ymin=-1.05, ymax=2.05,
        xlabel={$t$}, ylabel={$f(t)$},
        samples=400,
        domain=0:3,
        axis x line=middle, axis y line=middle,
        legend pos = north east,
        grid=both, minor tick num=1
      ]
        % Envelope
        \addplot[dashed, thick, gray, domain=0:3] {exp(-x)};
        \addlegendentry{$e^{-t}$}
        \addplot[dashed, thick, gray, domain=0:3] {-exp(-x)};
        \addlegendentry{$-e^{-t}$}

        % Real part
        \addplot[very thick, blue, domain=0:3] {exp(-x)*cos(deg(2*pi*x))};
        \addlegendentry{$\Re f(t)$}

        % Imag part
        \addplot[very thick, red, domain=0:3] {exp(-x)*sin(deg(2*pi*x))};
        \addlegendentry{$\Im f(t)$}

        % Vertical dashed lines for zeros
        \foreach \xx in {0,0.25,0.5,0.75,1,1.25,1.5,1.75,2,2.25,2.5,2.75} {
          \addplot[densely dashed, gray!40] coordinates {(\xx,-1.05) (\xx,1.05)};
        }
      \end{axis}
    \end{tikzpicture}


  \item Purely real / purely imaginary:
    \[
    \Im f(t)=0 \iff \sin(2\pi t)=0 \iff t=\frac{n}{2},\quad n\in\mathbb{Z}\quad(\text{purely real}).
    \]
    \[
    \Re f(t)=0 \iff \cos(2\pi t)=0 \iff t=\frac14+\frac{n}{2},\quad n\in\mathbb{Z}\quad(\text{purely imaginary}).
    \]
\end{enumerate}
}}
\end{center}


\subsection*{Question Q2}
Unforced oscillator:
\[
m\ddot x + c\dot x + kx = 0,\qquad m=4,\ k=9,\ c\ge0.
\]
\begin{enumerate}[(i)]
  \item Find characteristic roots.
  \item Identify ranges of \(c\) for (a) decaying without oscillation, (b) decaying with oscillation.
  \item Give one representative solution formula for each range.
\end{enumerate}

\begin{center}\setlength{\fboxsep}{8pt}\fcolorbox{yellow!20}{yellow!20}{\parbox{0.95\linewidth}{
\textbf{Solution (minimal).}

Characteristic equation:
\[
4r^2 + c r + 9 = 0 \quad\Rightarrow\quad r = \frac{-c \pm \sqrt{c^2 - 144}}{8}.
\]
Discriminant \(\Delta = c^2-144\).

Ranges:
\begin{itemize}
  \item \(\Delta>0\) (real distinct roots) $\iff c>12$. \(\Rightarrow\) decaying, no oscillation (overdamped).
  \item \(\Delta=0\) (double root) $\iff c=12$. critically damped (decaying, no oscillation).
  \item \(\Delta<0\) (complex roots) $\iff 0\le c<12$. \(\Rightarrow\) decaying with oscillation (underdamped).
\end{itemize}

Representative solutions:

\begin{enumerate}[(a)]
  \item Choose \(c=16>12\). Then
  \[
  r=\frac{-16\pm\sqrt{256-144}}{8}=\frac{-16\pm\sqrt{112}}{8}=\frac{-16\pm 4\sqrt7}{8}=\frac{-4\pm\sqrt7}{2}.
  \]
  General solution (overdamped):
  \[
    x(t)=C_1 e^{\frac{-4+\sqrt7}{2}t}+C_2 e^{\frac{-4-\sqrt7}{2}t}.
  \]

  \item Choose \(c=8\) (underdamped, \(0\le8<12\)). Then
  \[
  r=\frac{-8\pm\sqrt{64-144}}{8}=\frac{-8\pm i\sqrt{80}}{8}=-1\pm i\frac{\sqrt5}{2}.
  \]
  General solution (decaying oscillation):
  \[
    x(t)=e^{-t}\Big( A\cos\!\big(\tfrac{\sqrt5}{2}t\big)+B\sin\!\big(\tfrac{\sqrt5}{2}t\big)\Big).
  \]
\end{enumerate}
}}
\end{center}

\end{document}
